\section{Discussion}
\label{sec:discussion}

\subsection{RQ1: Workload and Scaling Sensitivity}

\subsubsection{Read/Write Composition}

The read/write-ratio experiment (Figure~\ref{fig:throughput-write-ratio}) shows that throughput does not vary linearly with workload composition. A first-order linear regression did not provide sufficient evidence for a linear relationship, whereas a second-order regression captured the observed curvature. The fitted responses showed a lower-throughput region for mixed read/write workloads and higher throughput towards the read- and write-dominated extremes.

The stronger curvature observed for EPaxos indicates that its throughput changed more substantially as the workload moved towards either extreme. This should primarily be interpreted as greater workload sensitivity rather than as evidence of uniformly higher efficiency. The two extremes expose different processing characteristics. In read-dominated workloads, reads do not introduce the same replicated state mutation and log-commit processing associated with writes. In write-dominated workloads, the evaluated EPaxos implementation can distribute proposal processing across replicas, whereas the evaluated Raft configuration directs requests through a single elected leader.

This distinction is particularly relevant to write scalability. In the evaluated Raft configuration, write requests are directed to the elected leader and submitted through the Raft state machine. Consequently, ordering, replication, and associated request-processing work remain concentrated at that replica. Adding replicas therefore does not directly distribute the write-processing path across the cluster. Replica addition can increase the number of available followers for replication and fault tolerance, but it does not by itself create additional leader-side request-processing capacity.

EPaxos has a different processing structure because proposals can be initiated at multiple replicas. Under sufficiently concurrent write workloads, this provides the opportunity to distribute proposal-processing work across the replica set. The higher recovery of EPaxos towards the write-dominated end of the workload sweep is therefore consistent with this architectural difference. However, the curvature alone does not establish a causal relationship between workload composition and a particular internal processing cost; it shows how the measured throughput responds to the composition under the evaluated implementation and environment.

Overall, the read/write experiment demonstrates why a single workload ratio can obscure important differences between consensus implementations. The same protocol may exhibit different relative behaviour depending on whether the workload exposes read processing, replicated writes, or distributed proposal processing.

\subsubsection{Replica Scaling}

Increasing the replica count reduced throughput for both protocols (Figure~\ref{fig:throughput-scaling}, Table~\ref{tab:scaling}) and increased median latency only for Raft (Figure~\ref{fig:latency-replica-scaling}, Table~\ref{tab:latency}). EPaxos nevertheless maintained a higher mean throughput than Raft at every tested replica count, with the relative difference ranging from approximately $1.29\times$ to $1.75\times$. The largest mean difference occurred at five replicas. However, the 95\,\% confidence intervals of the means overlapped at seven and nine replicas, so the magnitude of the difference at those cluster sizes should not be treated as conclusive.

The reduction in throughput with additional replicas is therefore not specific to leader-based consensus. Both protocols incur additional communication and coordination work as the replica set grows. In EPaxos, larger replica sets can additionally increase the amount of dependency and coordination information associated with proposal processing. The evaluated EPaxos implementation was also modified to support the dependency-set sizes required for seven- and nine-replica configurations, so results at those sizes include implementation-specific effects in addition to the protocol-level cost of larger clusters.

The Raft results expose a different scaling constraint. Increasing the number of followers increases the amount of replication work associated with a write, but the request-processing path remains centred on the elected leader. Consequently, replica addition does not horizontally distribute the primary write-processing workload. This provides a structural explanation for why adding replicas does not produce additional write throughput in the same way that adding independent request-processing endpoints could.

The scaling experiment therefore separates two effects that can otherwise be conflated. Increasing the replica count provides additional redundancy for both approaches, but it does not necessarily provide proportional additional processing capacity. In particular, the evaluated Raft configuration retains a centralised write path even as the cluster grows.

\subsubsection{Client Concurrency}

Client concurrency produced a clearer difference in throughput scaling (Figure~\ref{fig:throughput-concurrency}; the corresponding median latency is shown in Figure~\ref{fig:latency-concurrency}). At low concurrency, the measured throughputs were comparatively close, whereas increasing concurrency produced a substantially larger throughput increase for EPaxos. At the highest tested concurrency, the difference between the protocols was considerably larger than at low concurrency.

This behaviour is consistent with the distribution of request-processing work. In the evaluated Raft configuration, requests are directed to the elected leader. Increasing client concurrency therefore increases the amount of work presented to the same request-processing endpoint. EPaxos, in contrast, permits proposals to be initiated at multiple replicas, allowing additional concurrency to expose processing capacity distributed across the replica set.

The result suggests that centralised request processing becomes more visible as offered concurrency increases. The experiment does not establish that EPaxos universally scales better with concurrency, because the measured response also depends on implementation and workload characteristics. Rather, under the evaluated configuration, the leader-centred request path imposed a stronger throughput constraint at high concurrency.

The concurrency result also complements the read/write experiment. The stronger throughput recovery of EPaxos towards write-dominated workloads is consistent with distributed proposal processing, while the concurrency sweep identifies a workload condition under which that structural difference becomes more apparent.

\subsubsection{Conflict Rate}

The conflict-ratio experiment (Figure~\ref{fig:throughput-conflict-ratio}) did not reveal a clear monotonic relationship between conflict ratio and throughput. Intermediate measurements varied non-monotonically for both protocols, and no simple linear trend described the complete set of observations. Therefore, the results do not support interpreting every incremental increase in conflict as producing a corresponding throughput reduction.

A clearer difference appeared when comparing the extreme conflict conditions. Both protocols showed a throughput reduction between the lowest- and highest-conflict workloads, with the reduction being more pronounced for EPaxos. This indicates sensitivity to highly conflicting workloads without establishing a deterministic throughput function of conflict ratio.

The reduction observed for EPaxos at high conflict is consistent with the additional coordination required when its dependency relationships become more complex and commands are less likely to remain independent. In particular, the observation is compatible with the slow-path behaviour of EPaxos. However, the present measurements do not directly instrument the internal execution path sufficiently to establish that mechanism as the cause of the measured reduction. Additional instrumentation under controlled network and scheduling conditions would be required for such a causal analysis.

The conflict results therefore provide stronger evidence for sensitivity at the extremes than for a smooth relationship across the entire conflict range. This distinction is important because local scheduling and container-level variability can obscure relatively small changes between neighbouring conflict levels.

\subsection{RQ2: Failure Behaviour}

The failure experiments (Table~\ref{tab:failure}, Figure~\ref{fig:failure-recovery}) demonstrate a distinct difference between leader failure and non-leader replica failure. Killing the Raft leader produced a substantial availability gap, whereas killing a Raft follower caused little or no observable interruption in the measured workload. EPaxos replica failures also produced only short interruptions compared with the Raft leader-failure case. The Raft leader-kill experiments recorded availability gaps of 1.81--4.74\,s and 2{,}169--4{,}133 failed requests, while follower-kill and EPaxos replica-kill experiments recorded gaps below 0.1\,s with at most 106 failed requests.

The Raft result follows directly from the request path used in the experiment. The benchmark sends requests to the elected leader, so terminating that replica removes the endpoint through which requests are currently processed. The surviving replicas must subsequently complete leader detection and election before normal request processing can resume. The observed interruption therefore depends not only on the consensus protocol but also on election timing, failure detection, and client reconnection.

The election experiment (Figure~\ref{fig:election-failure}) further shows that the disruption is related to the difficulty of completing leader election. The measured number of failed election attempts was strongly associated with the observed availability gap ($r=0.90$, $n=60$). The association remained significant for recovery time measured after the isolation window had ended ($r=0.41$, $p=0.001$), indicating that election difficulty was not merely a consequence of the duration of the injected isolation period.

The follower-failure result provides an important contrast. When a Raft follower is removed, the elected leader remains available and the remaining replicas retain the quorum required by the evaluated three-replica configuration. The normal request path can therefore continue without the leader replacement step observed after leader failure.

EPaxos does not rely on a single leader endpoint for normal proposal placement. Consequently, failure of one replica does not require election of a replacement leader before other replicas can continue processing proposals. The short interruption observed after an EPaxos replica failure is consistent with this distributed proposal-processing structure.

These observations do not imply that leaderless consensus eliminates failure costs. A failed replica can still reduce available resources, affect quorum conditions, and introduce additional coordination. Rather, the experiment isolates a particular failure-mode difference: losing the replica responsible for the normal request-processing path has a substantially different availability effect from losing a non-leader replica in the evaluated systems.

The failure results therefore add a dimension that steady-state throughput cannot capture. A centralised leader can simplify request routing and provide a well-defined ordering point, but the same concentration creates a specific failure dependency. Conversely, distributing proposal processing removes that single request-processing endpoint, while retaining other costs associated with replica failure and coordination.

\subsection{RQ3: Implementation Sensitivity}

The implementation-sensitivity experiments show that measured consensus performance depends substantially on the implementation surrounding the protocol. Changes to persistence behaviour, request-processing paths, adapter-level processing, and communication conditions produced performance changes that were large relative to some of the baseline differences between the evaluated protocol implementations.

The clearest example was the Raft implementation. The initial implementation incurred substantial overhead from synchronous persistence in the request path. Removing this persistence bottleneck produced a large throughput increase, while allowing more asynchronous request processing produced further improvements and substantially reduced latency in the evaluated configurations. These changes altered measured performance without changing the underlying Raft consensus algorithm.

This observation is important for interpreting absolute throughput comparisons. A benchmark does not measure an abstract consensus algorithm in isolation; it measures a complete implementation stack consisting of consensus processing, state-machine execution, persistence, networking, serialisation, scheduling, and adapter behaviour. Consequently, a lower throughput measurement for one implementation cannot by itself be attributed to a fundamental limitation of the underlying consensus algorithm.

The communication-cost experiment (Figure~\ref{fig:throughput-commcost}) reinforces this point. The four evaluated implementations exhibited markedly different sensitivity to injected one-way latency. At 10\,ms of one-way cost, HashiCorp Raft lost approximately 72\,\% of its throughput, nvb EPaxos approximately 81\,\%, and the original EPaxos implementation approximately 96\,\%. The etcd Raft implementation showed an even sharper reduction, losing approximately 94\,\% of its throughput at only 1\,ms of injected cost. The corresponding latency measurements also diverged substantially across implementations (Figure~\ref{fig:latency-commcost}).

These differences cannot be explained solely by the distinction between leader-based and leaderless ordering, because implementations of the same protocol family exhibited substantially different responses to the same injected network condition. The result instead demonstrates that communication behaviour is mediated by implementation choices such as batching, message scheduling, transport handling, persistence, and request processing.

Jitter produced a comparatively smaller effect than the fixed communication cost in the tested configuration (Figure~\ref{fig:throughput-commcost-jitter}). At a 5\,ms one-way cost, increasing jitter to 100\,\% reduced throughput by approximately 24--36\,\% across the evaluated implementations. This suggests that the fixed communication delay was a more dominant factor than the tested variation in delay at that operating point. However, this observation should be interpreted only within the tested range and environment.

The implementation experiments therefore motivate a distinction between two levels of comparison. Absolute throughput and latency describe the behaviour of the particular implementation and environment being measured. Relative changes across controlled workload, concurrency, failure, and communication conditions provide evidence about how that implementation responds to changing conditions. The latter can remain informative even when absolute performance changes substantially between implementations.

At the same time, implementation sensitivity does not make protocol-level structure irrelevant. The workload and failure experiments show behaviour consistent with structural differences between the evaluated approaches: centralised request processing and leader dependence in the Raft configuration, versus distributed proposal placement and the absence of a single leader endpoint in EPaxos. The results indicate that these structural properties are observed through an implementation layer that can amplify, reduce, or obscure their performance effects.

Overall, the experiments support a cautious interpretation of consensus benchmarks. Absolute performance should be treated as a property of the evaluated implementation and experimental environment, while claims about protocol behaviour should rely on changes observed under controlled variations. This distinction is especially important when comparing implementations with different persistence, networking, scheduling, and execution architectures.

\subsection{Overall Interpretation}

Taken together, the three research questions show that the cost of centralised leadership is not represented by a single throughput difference. It appears in several interacting dimensions of the evaluated system.

First, the leader-centred request path limits how additional replicas can contribute to write-processing capacity in the evaluated Raft configuration. Replica addition increases the number of replicas participating in replication but does not move the primary request-processing path away from the leader. This constraint becomes more visible as client concurrency increases and is also reflected in the stronger recovery of EPaxos towards write-dominated workloads.

Second, leadership introduces a specific failure dependency. The Raft leader is both an ordering point and the endpoint used by the benchmark for normal request processing. Its failure therefore interrupts the request path and requires a new leader before normal processing can resume. The corresponding follower-failure experiment shows that this interruption is not an inherent consequence of losing any replica; it is specifically associated with losing the leader in the evaluated architecture.

Third, the magnitude of these effects is implementation-dependent. Persistence, request scheduling, adapter behaviour, and communication handling can substantially change the measured throughput and latency without changing the underlying consensus algorithm. The observed results should therefore not be interpreted as establishing a universal performance ordering between Raft and EPaxos.

The combined evidence instead supports a more specific conclusion: under the evaluated workloads and implementations, distributing proposal processing can reduce the dependence of request processing on a single replica, particularly under high concurrency and write-heavy workloads, while centralised leadership retains a simpler request-processing structure at the cost of concentrating part of the processing and failure path at the leader. The practical significance of this trade-off depends on workload composition, cluster size, communication conditions, failure requirements, and implementation design.
